% S1--S13 模型变体定义。
% 若使用下方表格中的 \toprule、\midrule 和 \bottomrule，请在主文件导言区加入：
% \usepackage{booktabs}

\subsection{S1--S13 模型变体定义}

为系统研究 tied embedding/unembedding 矩阵的结构化松绑方式及其对语言模型训练的影响，
本文在相同的 Transformer 主干上构造十三种 embedding-head 变体，记为 S1--S13。
所有变体共享相同的模型主干、优化器、训练数据与训练超参数，仅在输入 embedding
与输出 projection 的有效参数化方式上不同。

本文将“松绑”作为总称，表示在标准 tied 结构 $E_{\mathrm{eff}}=U_{\mathrm{eff}}=W$
之外引入额外自由度。下文中 LoRA 仅指词表维上的加性低秩增量
$AB$；隐藏维上的右乘低秩变换统一称为“乘性低秩松绑”，避免把所有低秩结构都混称为
LoRA。

设词表大小为 $V$，隐藏维度为 $d$。令共享的 tied embedding/unembedding 矩阵为
\begin{equation}
    W \in \mathbb{R}^{V \times d}.
\end{equation}
对于输入 token $x_t \in \{1,\dots,V\}$，其在共享矩阵中的 lookup 向量为
\begin{equation}
    W[x_t] \in \mathbb{R}^{d}.
\end{equation}
设 Transformer 最后一层输出 hidden state 为
\begin{equation}
    h_t \in \mathbb{R}^{d}.
\end{equation}
标准 tied 输出 logits 为
\begin{equation}
    z_t = W h_t + b,
\end{equation}
其中 $b \in \mathbb{R}^{V}$ 为输出偏置项。

本文首先将 S1--S7 统一理解为对 tied 矩阵 $W$ 的不同单侧松绑方式，并进一步在 S8--S10
中考察 untied 输入 embedding 矩阵 $E$ 上的三种输入侧松绑消融；S11--S13 则补充输出侧
平移、输入侧平移与乘性组合、以及输入侧加性与乘性组合三种扩展。具体而言，不同变体的核心
区别在于输入侧实际使用的有效矩阵 $E_{\mathrm{eff}}$ 和输出侧实际使用的有效矩阵
$U_{\mathrm{eff}}$：
\begin{equation}
    e_t = E_{\mathrm{eff}}[x_t], \qquad
    z_t = U_{\mathrm{eff}} h_t + b.
\end{equation}

\paragraph{S1: 标准 tied，不松绑。}
S1 是最基础的 weight tying 结构。输入侧和输出侧完全使用同一个矩阵：
\begin{equation}
    E_{\mathrm{eff}} = W, \qquad
    U_{\mathrm{eff}} = W.
\end{equation}
因此输入表示为
\begin{equation}
    e_t = W[x_t],
\end{equation}
输出 logits 为
\begin{equation}
    z_t = W h_t + b.
\end{equation}

也就是说，S1 没有任何额外自由度，输入 embedding 和输出 unembedding 被强制完全相同。
该模型作为标准 tied baseline，用于衡量其他松绑方式带来的结构变化。

\paragraph{S2: Full untied，完全松绑。}
S2 是最彻底的松绑方式。输入侧和输出侧各自学习一个完整矩阵：
\begin{equation}
    E_{\mathrm{eff}} = E, \qquad
    U_{\mathrm{eff}} = U,
\end{equation}
其中
\begin{equation}
    E \in \mathbb{R}^{V \times d}, \qquad
    U \in \mathbb{R}^{V \times d}, \qquad
    E \neq U.
\end{equation}
输入表示为
\begin{equation}
    e_t = E[x_t],
\end{equation}
输出 logits 为
\begin{equation}
    z_t = U h_t + b.
\end{equation}

S2 的自由度最大，但参数成本也最高。相较于标准 tied 结构，它额外引入了一整套
$V \times d$ 规模的矩阵参数，用于完全解除输入 embedding 与输出 unembedding 之间的共享约束。

\paragraph{S3: $\beta$ 向量松绑，输入侧 rank-1 全局平移松绑。}
S3 保留共享矩阵 $W$，但在输入侧加入一个可学习向量：
\begin{equation}
    \beta \in \mathbb{R}^{d}.
\end{equation}
输入侧有效矩阵可以写成
\begin{equation}
    E_{\mathrm{eff}} = W - \mathbf{1}\beta^\top,
\end{equation}
其中
\begin{equation}
    \mathbf{1} \in \mathbb{R}^{V}
\end{equation}
是全 1 向量。因此，$\mathbf{1}\beta^\top \in \mathbb{R}^{V \times d}$，
且该矩阵的每一行都是同一个 $\beta^\top$。

对单个 token 来说，输入表示为
\begin{equation}
    e_t = W[x_t] - \beta.
\end{equation}
输出侧不变：
\begin{equation}
    U_{\mathrm{eff}} = W,
\end{equation}
因此
\begin{equation}
    z_t = W h_t + b.
\end{equation}

该模型可以理解为一种非常特殊的加性低秩松绑：
\begin{equation}
    E_{\mathrm{eff}} = W + \Delta W,
\end{equation}
其中
\begin{equation}
    \Delta W = -\mathbf{1}\beta^\top.
\end{equation}
由于 $\mathbf{1}\beta^\top$ 是一个 rank-1 矩阵，S3 本质上是 rank-1 的全局平移松绑。
它不是给每个 token 单独学习修正，而是让所有 token 共享同一个平移方向。

\paragraph{S4: 输入侧词表维加性低秩松绑（LoRA 形式）。}
S4 在输入 embedding 矩阵上加入词表维加性低秩增量。输入侧有效矩阵为
\begin{equation}
    E_{\mathrm{eff}} = W + \Delta W_{\mathrm{in}},
\end{equation}
其中
\begin{equation}
    \Delta W_{\mathrm{in}} = A_{\mathrm{in}}B_{\mathrm{in}},
\end{equation}
并且
\begin{equation}
    A_{\mathrm{in}} \in \mathbb{R}^{V \times r}, \qquad
    B_{\mathrm{in}} \in \mathbb{R}^{r \times d}, \qquad
    r \ll d.
\end{equation}
因此
\begin{equation}
    E_{\mathrm{eff}} = W + A_{\mathrm{in}}B_{\mathrm{in}},
\end{equation}
输入表示为
\begin{equation}
    e_t = \left(W + A_{\mathrm{in}}B_{\mathrm{in}}\right)[x_t].
\end{equation}
输出侧仍然使用原始共享矩阵：
\begin{equation}
    U_{\mathrm{eff}} = W,
\end{equation}
即
\begin{equation}
    z_t = W h_t + b.
\end{equation}

S4 的含义是只松绑输入 embedding，不松绑输出 unembedding。该方式比 full untied
参数量更低，但由于 $A_{\mathrm{in}} \in \mathbb{R}^{V \times r}$，额外参数量仍然包含
\begin{equation}
    Vr + rd.
\end{equation}
当词表大小 $V$ 很大时，$Vr$ 这一项仍然会带来明显开销。

\paragraph{S5: 输出侧词表维加性低秩松绑（LoRA 形式）。}
S5 与 S4 对称，但加性低秩增量作用在输出 projection 上。输入侧保持 tied：
\begin{equation}
    E_{\mathrm{eff}} = W,
\end{equation}
因此
\begin{equation}
    e_t = W[x_t].
\end{equation}
输出侧有效矩阵为
\begin{equation}
    U_{\mathrm{eff}} = W + \Delta W_{\mathrm{out}},
\end{equation}
其中
\begin{equation}
    \Delta W_{\mathrm{out}} = A_{\mathrm{out}}B_{\mathrm{out}},
\end{equation}
并且
\begin{equation}
    A_{\mathrm{out}} \in \mathbb{R}^{V \times r}, \qquad
    B_{\mathrm{out}} \in \mathbb{R}^{r \times d}, \qquad
    r \ll d.
\end{equation}
因此
\begin{equation}
    U_{\mathrm{eff}} = W + A_{\mathrm{out}}B_{\mathrm{out}},
\end{equation}
输出 logits 为
\begin{equation}
    z_t =
    \left(W + A_{\mathrm{out}}B_{\mathrm{out}}\right)h_t + b.
\end{equation}

S5 的含义是输入侧仍然 tied，输出侧通过词表维加性低秩增量获得额外自由度。S4 和 S5
构成一组输入侧加性低秩与输出侧加性低秩的结构对照。

\paragraph{S6: 输入侧乘性低秩松绑。}
S6 是输入侧的乘性低秩松绑方式。由于语言模型中通常有 $V \gg d$，直接在
$V \times d$ 的 embedding 矩阵上加入词表维加性低秩增量会引入与词表大小 $V$ 相关的参数。
因此，S6 不在词表维度上引入低秩矩阵，而是在隐藏维度上右乘一个低秩变换。

输入侧有效矩阵写成
\begin{equation}
    E_{\mathrm{eff}} = W M_{\mathrm{in}},
\end{equation}
其中
\begin{equation}
    M_{\mathrm{in}} \in \mathbb{R}^{d \times d}.
\end{equation}
为了保持接近原始 tied embedding，将 $M_{\mathrm{in}}$ 写成单位矩阵加低秩修正：
\begin{equation}
    M_{\mathrm{in}} = I_d + \Delta M_{\mathrm{in}},
\end{equation}
其中
\begin{equation}
    \Delta M_{\mathrm{in}} = P_{\mathrm{in}}Q_{\mathrm{in}},
\end{equation}
并且
\begin{equation}
    P_{\mathrm{in}} \in \mathbb{R}^{d \times r}, \qquad
    Q_{\mathrm{in}} \in \mathbb{R}^{r \times d}, \qquad
    r \ll d.
\end{equation}
所以
\begin{equation}
    E_{\mathrm{eff}} =
    W(I_d + P_{\mathrm{in}}Q_{\mathrm{in}}).
\end{equation}
对单个 token，有
\begin{equation}
    e_t =
    W[x_t](I_d + P_{\mathrm{in}}Q_{\mathrm{in}}).
\end{equation}
输出侧不变：
\begin{equation}
    U_{\mathrm{eff}} = W,
\end{equation}
即
\begin{equation}
    z_t = W h_t + b.
\end{equation}

S6 的含义是不额外学习一个 $V \times r$ 的词表侧低秩因子，而是在 hidden dimension
上学习一个 $d \times d$ 的低秩变换。其额外参数量为
\begin{equation}
    dr + rd = 2dr,
\end{equation}
不随词表大小 $V$ 增长。需要强调的是，这里不是额外插入一个普通线性层，而是把输入
embedding 矩阵重新参数化为
\begin{equation}
    E_{\mathrm{eff}} =
    W(I_d + P_{\mathrm{in}}Q_{\mathrm{in}}).
\end{equation}
也就是说，模型实际使用的是一个被乘性松绑后的有效输入矩阵。

\paragraph{S7: 输出侧乘性低秩松绑。}
S7 是 S6 的输出侧版本。输入侧保持 tied：
\begin{equation}
    E_{\mathrm{eff}} = W,
\end{equation}
因此
\begin{equation}
    e_t = W[x_t].
\end{equation}
输出侧使用乘性松绑矩阵：
\begin{equation}
    U_{\mathrm{eff}} = W M_{\mathrm{out}},
\end{equation}
其中
\begin{equation}
    M_{\mathrm{out}} = I_d + \Delta M_{\mathrm{out}},
\end{equation}
\begin{equation}
    \Delta M_{\mathrm{out}} = P_{\mathrm{out}}Q_{\mathrm{out}},
\end{equation}
并且
\begin{equation}
    P_{\mathrm{out}} \in \mathbb{R}^{d \times r}, \qquad
    Q_{\mathrm{out}} \in \mathbb{R}^{r \times d}, \qquad
    r \ll d.
\end{equation}
因此
\begin{equation}
    U_{\mathrm{eff}} =
    W(I_d + P_{\mathrm{out}}Q_{\mathrm{out}}),
\end{equation}
输出 logits 为
\begin{equation}
    z_t =
    W(I_d + P_{\mathrm{out}}Q_{\mathrm{out}})h_t + b.
\end{equation}

S7 的含义是不在词表维度上给输出矩阵加加性低秩增量，而是在输出 projection 的隐藏维度方向
做乘性低秩松绑。S6 和 S7 构成一组输入侧乘性松绑与输出侧乘性松绑的结构对照。

\paragraph{S8: Untied Embedding + $\beta$，独立输入矩阵上的 rank-1 全局平移松绑。}
S8 以 S2 的 full untied 结构为基础，但进一步在独立输入 embedding 矩阵 $E$ 上加入
rank-1 全局平移松绑。输出侧仍然使用独立的输出矩阵 $U$。

输入侧有效矩阵为
\begin{equation}
    E_{\mathrm{eff}} = E - \mathbf{1}\beta_E^\top,
\end{equation}
其中
\begin{equation}
    \beta_E \in \mathbb{R}^{d}, \qquad
    \mathbf{1} \in \mathbb{R}^{V}.
\end{equation}
对单个 token，有
\begin{equation}
    e_t = E[x_t] - \beta_E.
\end{equation}
输出侧保持 untied：
\begin{equation}
    U_{\mathrm{eff}} = U,
\end{equation}
因此
\begin{equation}
    z_t = U h_t + b.
\end{equation}

S8 用于检验在输入输出已经完全 untied 的情况下，输入 embedding 层是否仍然受益于
rank-1 的全局平移校正。与 S3 不同，S8 的平移作用在独立输入矩阵 $E$ 上，而不是
共享矩阵 $W$ 上。

\paragraph{S9: Untied Embedding + Additive Low-Rank，独立输入矩阵上的词表维加性低秩松绑。}
S9 同样以 S2 的 full untied 结构为基础，但在独立输入 embedding 矩阵 $E$ 上加入
词表维加性低秩修正。输入侧有效矩阵为
\begin{equation}
    E_{\mathrm{eff}} = E + \Delta E_{\mathrm{in}},
\end{equation}
其中
\begin{equation}
    \Delta E_{\mathrm{in}} = A_{E}B_{E},
\end{equation}
并且
\begin{equation}
    A_{E} \in \mathbb{R}^{V \times r}, \qquad
    B_{E} \in \mathbb{R}^{r \times d}, \qquad
    r \ll d.
\end{equation}
因此
\begin{equation}
    e_t = \left(E + A_EB_E\right)[x_t].
\end{equation}
输出侧仍然为
\begin{equation}
    U_{\mathrm{eff}} = U,
\end{equation}
即
\begin{equation}
    z_t = U h_t + b.
\end{equation}

S9 用于考察词表维加性低秩修正是否能够在 full untied 的输入 embedding 上带来额外收益。
与 S4 不同，S9 的低秩加性修正不是对共享矩阵 $W$ 进行松绑，而是对已经独立学习的
输入矩阵 $E$ 进行进一步重参数化。

\paragraph{S10: Untied Embedding + Multiplicative Low-Rank，独立输入矩阵上的乘性低秩松绑。}
S10 是 S8--S9 的乘性版本。它不在词表维度上对 $E$ 加入 $V \times r$ 的低秩增量，
而是在隐藏维度上右乘一个低秩变换。输入侧有效矩阵为
\begin{equation}
    E_{\mathrm{eff}} = E M_E,
\end{equation}
其中
\begin{equation}
    M_E = I_d + \Delta M_E,
\end{equation}
\begin{equation}
    \Delta M_E = P_E Q_E,
\end{equation}
并且
\begin{equation}
    P_E \in \mathbb{R}^{d \times r}, \qquad
    Q_E \in \mathbb{R}^{r \times d}, \qquad
    r \ll d.
\end{equation}
因此
\begin{equation}
    E_{\mathrm{eff}} = E(I_d+P_EQ_E),
\end{equation}
对单个 token，有
\begin{equation}
    e_t = E[x_t](I_d+P_EQ_E).
\end{equation}
输出侧保持独立矩阵：
\begin{equation}
    U_{\mathrm{eff}} = U,
\end{equation}
即
\begin{equation}
    z_t = U h_t + b.
\end{equation}

S10 用于检验在 full untied 的输入 embedding 上，隐藏维度乘性松绑是否比词表维度加性低秩修正
更具参数效率。其额外参数量为 $2dr$，不随词表大小 $V$ 增长。与 S6 不同，S10 的乘性
低秩变换作用在独立输入矩阵 $E$ 上，而不是共享矩阵 $W$ 上。

\paragraph{S11: 输出侧隐藏维全局平移松绑。}
S11 保持输入侧 tied，不直接改变输出矩阵 $W$，而是在输出 projection 前对 hidden state
加入一个全局平移参数：
\begin{equation}
    \beta_{\mathrm{out}} \in \mathbb{R}^{d}.
\end{equation}
输入侧为
\begin{equation}
    E_{\mathrm{eff}} = W,
\end{equation}
输出 logits 写为
\begin{equation}
    z_t = W(h_t-\beta_{\mathrm{out}})+b.
\end{equation}

因此，S11 可以看作输出侧 hidden 维的全局平移松绑。它与 S3 对称：S3 在输入
embedding lookup 后做平移，S11 在输出 projection 前对 hidden state 做平移。需要注意，
这里的 $\beta_{\mathrm{out}}$ 不是 $V$ 维的 lm head bias；lm head bias 是所有新版实验中
单独控制的输出偏置项 $b$。

\paragraph{S12: 输入侧全局平移 + 乘性低秩组合。}
S12 将 S3 的输入侧全局平移与 S6 的输入侧乘性低秩松绑组合起来。给定原始 lookup
\begin{equation}
    e_t^{(0)} = W[x_t],
\end{equation}
实际输入表示为
\begin{equation}
    e_t = e_t^{(0)}-\beta + e_t^{(0)}P_{\mathrm{in}}Q_{\mathrm{in}}.
\end{equation}
等价地，输入侧有效矩阵为
\begin{equation}
    E_{\mathrm{eff}} =
    W(I_d+P_{\mathrm{in}}Q_{\mathrm{in}})-\mathbf{1}\beta^\top,
\end{equation}
输出侧仍然保持
\begin{equation}
    U_{\mathrm{eff}} = W.
\end{equation}

S12 的额外参数量为 $d+2dr$。它不是 full untied，也不引入词表维加性低秩因子，而是在
输入侧同时提供一个全局平移自由度和一个隐藏维低秩线性自由度。

\paragraph{S13: 输入侧词表维加性低秩 + 乘性低秩组合。}
S13 将 S4 的输入侧词表维加性低秩增量与 S6 的输入侧乘性低秩松绑组合起来。输入侧有效矩阵为
\begin{equation}
    E_{\mathrm{eff}} =
    W + A_{\mathrm{in}}B_{\mathrm{in}} + WP_{\mathrm{in}}Q_{\mathrm{in}},
\end{equation}
即对单个 token 有
\begin{equation}
    e_t =
    W[x_t] + \left(A_{\mathrm{in}}B_{\mathrm{in}}\right)[x_t]
    + W[x_t]P_{\mathrm{in}}Q_{\mathrm{in}}.
\end{equation}
输出侧仍然保持 tied：
\begin{equation}
    U_{\mathrm{eff}} = W.
\end{equation}

S13 用于检验词表维加性低秩自由度与隐藏维乘性低秩自由度是否互补。其额外参数量为
\begin{equation}
    Vr + rd + 2dr,
\end{equation}
其中第一项来自词表维加性低秩因子，后两项来自隐藏维乘性低秩因子。

\subsection{十三个模型的统一对比}

\begin{table}[htbp]
\centering
\caption{S1--S13 模型变体的有效输入矩阵与有效输出矩阵。}
\label{tab:s1_s13_variants}
\begin{tabular}{llll}
\toprule
模型 & 松绑方式 & $E_{\mathrm{eff}}$ & $U_{\mathrm{eff}}$ \\
\midrule
S1 & 标准 tied（不松绑） & $W$ & $W$ \\
S2 & full untied（完全松绑） & $E$ & $U$ \\
S3 & 输入侧 rank-1 全局平移松绑 & $W-\mathbf{1}\beta^\top$ & $W$ \\
S4 & 输入侧词表维加性低秩松绑 & $W+A_{\mathrm{in}}B_{\mathrm{in}}$ & $W$ \\
S5 & 输出侧词表维加性低秩松绑 & $W$ & $W+A_{\mathrm{out}}B_{\mathrm{out}}$ \\
S6 & 输入侧乘性低秩松绑 & $W(I_d+P_{\mathrm{in}}Q_{\mathrm{in}})$ & $W$ \\
S7 & 输出侧乘性低秩松绑 & $W$ & $W(I_d+P_{\mathrm{out}}Q_{\mathrm{out}})$ \\
S8 & untied 输入侧 rank-1 全局平移松绑 & $E-\mathbf{1}\beta_E^\top$ & $U$ \\
S9 & untied 输入侧词表维加性低秩松绑 & $E+A_EB_E$ & $U$ \\
S10 & untied 输入侧乘性低秩松绑 & $E(I_d+P_EQ_E)$ & $U$ \\
S11 & 输出侧隐藏维全局平移松绑 & $W$ & $W(h_t-\beta_{\mathrm{out}})+b$ \\
S12 & 输入侧全局平移 + 乘性低秩松绑 & $W(I_d+P_{\mathrm{in}}Q_{\mathrm{in}})-\mathbf{1}\beta^\top$ & $W$ \\
S13 & 输入侧加性低秩 + 乘性低秩松绑 & $W+A_{\mathrm{in}}B_{\mathrm{in}}+WP_{\mathrm{in}}Q_{\mathrm{in}}$ & $W$ \\
\bottomrule
\end{tabular}
\end{table}

\paragraph{初始化策略。}
为保证各变体间的公平比较，所有新增参数均以恒等映射或零初始化的方式引入：对于乘性低秩变换
$I_d+P_{\mathrm{in}}Q_{\mathrm{in}}$ 和 $I_d+P_{\mathrm{out}}Q_{\mathrm{out}}$，其中
$P$（即 \texttt{mul\_a}）以 $\mathcal{N}(0,0.02)$ 初始化，$Q$（即 \texttt{mul\_b}）以全零
初始化，使得训练起点 $M_{\mathrm{in}}=M_{\mathrm{out}}=I_d$；对于输入侧加性低秩增量
$A_{\mathrm{in}}B_{\mathrm{in}}$，$A_{\mathrm{in}}$（即 \texttt{lora\_a}）以
$\mathcal{N}(0,0.02)$ 初始化，$B_{\mathrm{in}}$（即 \texttt{lora\_b}）以全零初始化；对于
输出侧加性低秩增量 $A_{\mathrm{out}}B_{\mathrm{out}}$ 同理；$\beta$ 类参数以全零初始化。
因此，在训练的第 0 步，所有变体 S1--S13 的有效输入输出矩阵均退化为 $W$ 或 $E/U$，
确保训练起点的等价性。

\subsection{十三个模型的逻辑关系}

上述十三个模型可以分为七类。

\paragraph{第一类：基线模型。}
S1 为标准 tied（不松绑）的基线：
\begin{equation}
    E_{\mathrm{eff}} = W, \qquad
    U_{\mathrm{eff}} = W.
\end{equation}
S2 为 full untied（完全松绑）的上界参照：
\begin{equation}
    E_{\mathrm{eff}} = E, \qquad
    U_{\mathrm{eff}} = U.
\end{equation}

\paragraph{第二类：输入侧加性松绑。}
S3 和 S4 都只改变输入侧，输出侧仍然满足
\begin{equation}
    U_{\mathrm{eff}} = W.
\end{equation}
其中，S3 使用 rank-1 的全局平移松绑：
\begin{equation}
    E_{\mathrm{eff}} = W-\mathbf{1}\beta^\top,
\end{equation}
S4 使用输入侧词表维加性低秩松绑：
\begin{equation}
    E_{\mathrm{eff}} = W+A_{\mathrm{in}}B_{\mathrm{in}}.
\end{equation}

\paragraph{第三类：输出侧加性松绑。}
S5 和 S11 都只改变输出侧，输入侧仍然满足
\begin{equation}
    E_{\mathrm{eff}} = W,
\end{equation}
其中，S5 在输出矩阵上加入词表维加性低秩增量：
\begin{equation}
    U_{\mathrm{eff}} = W+A_{\mathrm{out}}B_{\mathrm{out}}.
\end{equation}
S11 则不改变 $W$ 本身，而是在输出 projection 前平移 hidden state：
\begin{equation}
    z_t = W(h_t-\beta_{\mathrm{out}})+b.
\end{equation}

\paragraph{第四类：乘性低秩松绑。}
S6 和 S7 将低秩自由度放在隐藏维度 $d$ 上，而不是词表维度 $V$ 上。
输入侧乘性松绑为
\begin{equation}
    \text{S6:}\quad
    E_{\mathrm{eff}} = W(I_d+P_{\mathrm{in}}Q_{\mathrm{in}}),
\end{equation}
输出侧乘性松绑为
\begin{equation}
    \text{S7:}\quad
    U_{\mathrm{eff}} = W(I_d+P_{\mathrm{out}}Q_{\mathrm{out}}).
\end{equation}
这两种方法的额外参数量均为 $2dr$，不随词表大小 $V$ 增长。

\paragraph{第五类：untied 输入 embedding 上的松绑消融。}
S8--S10 以 S2 的 full untied 结构为基础，保持输出侧为独立矩阵
\begin{equation}
    U_{\mathrm{eff}} = U,
\end{equation}
并只在独立输入 embedding 矩阵 $E$ 上加入额外松绑。三种形式分别为
\begin{equation}
    \text{S8:}\quad
    E_{\mathrm{eff}} = E-\mathbf{1}\beta_E^\top,
\end{equation}
\begin{equation}
    \text{S9:}\quad
    E_{\mathrm{eff}} = E+A_EB_E,
\end{equation}
以及
\begin{equation}
    \text{S10:}\quad
    E_{\mathrm{eff}} = E(I_d+P_EQ_E).
\end{equation}
这组三个模型用于判断本文讨论的三类输入侧松绑形式是否只在 tied 约束下有效，还是在
输入输出已经完全 untied 的情况下仍然具有额外作用。

\paragraph{第六类：输入侧平移与乘性低秩组合。}
S12 将 S3 与 S6 组合，在输入侧同时引入全局平移和隐藏维乘性低秩自由度：
\begin{equation}
    \text{S12:}\quad
    E_{\mathrm{eff}} =
    W(I_d+P_{\mathrm{in}}Q_{\mathrm{in}})-\mathbf{1}\beta^\top.
\end{equation}
该模型用于检验输入侧平移与乘性低秩变换是否互补。

\paragraph{第七类：输入侧加性低秩与乘性低秩组合。}
S13 将 S4 与 S6 组合，在输入侧同时引入词表维加性低秩增量和隐藏维乘性低秩变换：
\begin{equation}
    \text{S13:}\quad
    E_{\mathrm{eff}} =
    W+A_{\mathrm{in}}B_{\mathrm{in}}+WP_{\mathrm{in}}Q_{\mathrm{in}}.
\end{equation}
该模型用于判断两类低秩自由度是否提供互补收益。

\subsection{方法概述}

总体而言，本文将不同模型统一理解为对 embedding/unembedding matrix 的结构化松绑。
标准 tied 模型直接令输入矩阵和输出矩阵均等于 $W$；full untied 模型则完全解除该约束，
分别学习独立的输入矩阵 $E$ 和输出矩阵 $U$。在此基础上，我们进一步比较三类轻量松绑方式：
单向量平移松绑、词表维加性低秩松绑（LoRA 形式），以及隐藏维度上的乘性低秩松绑；同时，
我们在 untied 输入 embedding 矩阵 $E$ 上加入同样三种形式，作为输入层松绑的消融对照。
当前版本还包含 S11--S13：S11 是输出侧 hidden 平移，S12 是输入侧平移与乘性低秩组合，
S13 是输入侧词表维加性低秩与乘性低秩组合。

对于单向量平移松绑，我们令
\begin{equation}
    E_{\mathrm{eff}} = W-\mathbf{1}\beta^\top,
\end{equation}
其中 $\beta\in\mathbb{R}^{d}$，$\mathbf{1}\in\mathbb{R}^{V}$。该形式等价于对 $W$
施加一个 rank-1 的加性修正。

对于词表维加性低秩松绑，我们分别在输入侧或输出侧引入低秩加性矩阵：
\begin{equation}
    E_{\mathrm{eff}} = W+A_{\mathrm{in}}B_{\mathrm{in}},
\end{equation}
或
\begin{equation}
    U_{\mathrm{eff}} = W+A_{\mathrm{out}}B_{\mathrm{out}}.
\end{equation}
其中 $A\in\mathbb{R}^{V\times r}$，$B\in\mathbb{R}^{r\times d}$。
该方法直接在 $V\times d$ 的词表矩阵上施加低秩修正，也就是本文中保留 LoRA 称呼的情形。

考虑到语言模型中通常 $V\gg d$，我们进一步引入乘性松绑，将有效矩阵写为共享矩阵右乘
一个隐藏维度上的低秩变换：
\begin{equation}
    E_{\mathrm{eff}} = W(I_d+P_{\mathrm{in}}Q_{\mathrm{in}}),
\end{equation}
或
\begin{equation}
    U_{\mathrm{eff}} = W(I_d+P_{\mathrm{out}}Q_{\mathrm{out}}).
\end{equation}
其中 $P\in\mathbb{R}^{d\times r}$，$Q\in\mathbb{R}^{r\times d}$。
该设计只在隐藏维度上引入 $2dr$ 个额外参数，不随词表大小 $V$ 增长。
从计算等价性上看，该乘性形式等同于在 embedding 特征后施加一个带残差结构的线性变换；
但在实现与推理时，我们直接对输入或输出 embedding matrix 本身进行重参数化，从而避免
修改模型主干结构，并减少额外算子带来的推理开销。

对于 untied 输入 embedding 的消融模型，我们将上述三类松绑形式从共享矩阵 $W$ 迁移到
独立输入矩阵 $E$ 上，即以下三种形式：
\begin{equation}
    E_{\mathrm{eff}} = E-\mathbf{1}\beta_E^\top,
\end{equation}
或者
\begin{equation}
    E_{\mathrm{eff}} = E+A_EB_E,
\end{equation}
或
\begin{equation}
    E_{\mathrm{eff}} = E(I_d+P_EQ_E),
\end{equation}
同时保持输出侧 $U_{\mathrm{eff}}=U$。该设置用于区分改进是否来自 tied 约束的缓解，
还是来自输入 embedding 层本身的额外重参数化能力。

对于组合模型，S12 在输入侧同时加入全局平移和隐藏维乘性低秩变换：
\begin{equation}
    E_{\mathrm{eff}} =
    W(I_d+P_{\mathrm{in}}Q_{\mathrm{in}})-\mathbf{1}\beta^\top,
\end{equation}
S13 则在输入侧同时加入词表维加性低秩增量和隐藏维乘性低秩变换：
\begin{equation}
    E_{\mathrm{eff}} =
    W+A_{\mathrm{in}}B_{\mathrm{in}}+WP_{\mathrm{in}}Q_{\mathrm{in}}.
\end{equation}
这两个组合模型保持输出侧 tied，用于检验不同输入侧自由度之间是否存在互补收益。
